% STATUS: written 2026-11-04 (Lu and Jing 2010 read in full, pp. 89--99; Benkovic 2015 read in full, pp. 141--165)
% STATUS: extended 2026-11-06 (Zhang and Wei 2017 read in full from the library PDF)
% Sources read in full: Lu and Jing 2010; Benkovic 2015; Zhang and Wei 2017 (pp. 1--3, 6).
\section{Lie derivations of specific algebras}\label{sec:lie-type:lie-on-algebras}

The sections so far moved through general classes of rings and algebras. On a
concrete algebra the question has a concrete answer: does the standard form
hold there, with what remainder, and what does the remainder look like. The two
answers this section records are the algebra of bounded operators on a Banach
space and the algebras built around an idempotent.

Lu and Jing settle the question of Martindale for $B(X)$, the algebra of all
bounded linear operators on a complex Banach space $X$ of dimension greater
than $2$ \cite[p.~91]{LuJing2010}. An additive map $\tau\colon B(X)\to B(X)$
is called \emph{centralized} when $\tau([A,B])\in\mathbb{C}I$ for all
$A,B\in B(X)$ \cite[p.~91]{LuJing2010}.

\begin{theorem}[{\cite[Theorem~2.1, p.~92]{LuJing2010}}]\label{thm:lie:lujing-bx}
Let $X$ be a Banach space with $\dim X>2$, and $\delta$ be a Lie derivation on
$B(X)$. Then there exists an operator $T\in B(X)$, and a linear map
$\tau\colon B(X)\to\mathbb{C}I$ vanishing on each commutator, such that
$\delta(A)=AT-TA+\tau(A)$ for all $A\in B(X)$.
\end{theorem}

Two points mark the answer.
Every Lie derivation of $B(X)$ is inner up to the
scalar remainder, and the remainder is finer than the general theory asks.
It need not be given by a linear functional $A\mapsto f(A)I$ on the subspace
spanned by the commutators \cite[p.~91]{LuJing2010}.
With the second point Lu and Jing answer in the negative a question of Brešar
\cite[p.~92]{LuJing2010}.

Benkovič gives the idempotent answer for linear maps over a commutative ring
$R$. The algebra $A$ is unital, $2$-torsion free, with an idempotent $e\ne0,1$
such that $exe\cdot eA(1-e)=\{0\}=(1-e)Ae\cdot exe$ implies $exe=0$ and
$eA(1-e)\cdot(1-e)x(1-e)=\{0\}=(1-e)x(1-e)\cdot(1-e)Ae$ implies
$(1-e)x(1-e)=0$ \cite[p.~141]{Benkovic2015}. A Jordan derivation
$\delta\colon A\to A$ is called \emph{singular} for the Peirce decomposition
when $\delta(eAe+(1-e)A(1-e))=0$, $\delta(eA(1-e))\subseteq(1-e)Ae$ and
$\delta((1-e)Ae)\subseteq eA(1-e)$ \cite[p.~144]{Benkovic2015}. Such a map is
a Lie triple derivation that no derivation plus central map can produce, so it
enters the final form \cite[p.~144]{Benkovic2015}.

\begin{theorem}[{\cite[Theorem~5.1, p.~159]{Benkovic2015}}]\label{thm:lie:benkovic-triple}
Let $A$ be a unital $2$-torsion free algebra with a nontrivial idempotent $e$
satisfying the annihilation condition above. Assume one of the conditions
\textup{(A)}, \textup{(B)}, \textup{(C)}, \textup{(D)}, on the corners $eAe$
and $(1-e)A(1-e)$: generation by idempotents and commutators, or the equality
of the corner centers with $Z(A)e$ and $Z(A)(1-e)$, and assume one of the
conditions \textup{(1)} to \textup{(4)}, the absence of central ideals in a
corner or a prescribed form of $Z(A)$. Then any Lie triple derivation
$D\colon A\to A$ is of the form $D=d+\delta+\gamma$, where $d\colon A\to A$
is a derivation, $\delta\colon A\to A$ is a singular Jordan derivation and
$\gamma\colon A\to Z(A)$ is a linear map that vanishes on $[[A,A],A]$
\cite[p.~159]{Benkovic2015}.
\end{theorem}

For Lie derivations the singular part disappears.

\begin{theorem}[{\cite[Theorem~5.3, p.~160]{Benkovic2015}}]\label{thm:lie:benkovic-lie}
Under the analogous hypotheses, a Lie derivation $D\colon A\to A$ is of the
form $D=d+\gamma$, where $d\colon A\to A$ is a derivation and
$\gamma\colon A\to Z(A)$ is a linear map that vanishes on $[A,A]$
\cite[p.~160]{Benkovic2015}.
\end{theorem}

The applications cover the standard families \cite[pp.~161--163]{Benkovic2015}.

\begin{corollary}[{\cite[Corollary~5.5, p.~161]{Benkovic2015}}]\label{cor:lie:benkovic-matrix}
Let $A=M_{n}(A)$, $n\ge3$, be a matrix algebra over a $2$-torsion free unital
algebra $A$. Then every Lie triple derivation $D\colon A\to A$ is of the form
$D=d+\gamma$, where $d$ is a derivation and $\gamma\colon A\to Z(A)$ vanishes
on $[[A,A],A]$.
\end{corollary}

The second family is the nest algebras.

\begin{corollary}[{\cite[Corollary~5.12, p.~163]{Benkovic2015}}]\label{cor:lie:benkovic-nest}
Let $\mathcal{N}$ be a nest of a complex Hilbert space $H$, $\dim H\ge2$. Then
every Lie triple derivation $D$ of the nest algebra $\mathcal{T}(\mathcal{N})$
is of the form $D=d+\gamma$, where $d$ is a derivation and
$\gamma\colon\mathcal{T}(\mathcal{N})\to\mathbb{C}I$ vanishes on
$[[\mathcal{A},\mathcal{A}],\mathcal{A}]$. In particular, every Lie derivation
of a nest algebra has the standard form.
\end{corollary}

The dimension $2$ case can fail.
Benkovič recalls an example from his earlier
work with Eremita of a triangular algebra $\mathcal{T}_{2}(A)$ over a
$\mathbb{Z}_{2}$-graded algebra.
On it, a Lie triple derivation is not the
sum of a derivation and a central linear mapping vanishing on the second
commutator \cite[p.~163]{Benkovic2015}. On the same algebra every Lie
derivation is still of the standard form \cite[p.~164]{Benkovic2015}.

On a nest algebra of a Banach space the same conclusion holds for every
Lie-type derivation, with no dimension bound
\cite[Theorem~3.1, p.~6]{ZhangWei2017Nest}. Zhang and Wei work with the
$(n-1)$-th commutator $p_n(x_1,\dots,x_n)=[\cdots[[x_1,x_2],x_3],\dots,x_n]$ and
call a linear map $L$ a Lie $n$-derivation when the Leibniz rule holds for
$p_n$ \cite[p.~2]{ZhangWei2017Nest}. Their theorem reads as follows
\cite[Theorem~3.1, p.~6]{ZhangWei2017Nest}.

\begin{theorem}[{\cite[Theorem~3.1, p.~6]{ZhangWei2017Nest}}]\label{thm:lie:zhang-wei-nest}
Let $X$ be a Banach space, $\mathcal{N}$ a non-trivial nest on $X$ and
$\operatorname{Alg}\mathcal{N}$ the nest algebra associated with $\mathcal{N}$.
Then $L\colon\operatorname{Alg}\mathcal{N}\to B(X)$ is a Lie $n$-derivation if
and only if there exist a derivation $D\colon\operatorname{Alg}\mathcal{N}\to B(X)$
and a linear map $H\colon\operatorname{Alg}\mathcal{N}\to\mathbb{C}I$ that
vanishes on each $(n-1)$-th commutator, such that $L=D+H$. If $L$ is continuous,
then $D$ and $H$ are continuous as well.
\end{theorem}

So every Lie-type derivation of a nest algebra on a Banach space has the
standard form \cite[p.~3]{ZhangWei2017Nest}. The authors prove the theorem for
the algebra itself and for subalgebras that meet six conditions on their
rank-one operators \cite[pp.~6--7]{ZhangWei2017Nest}. The nest algebra meets
them \cite[p.~7]{ZhangWei2017Nest}.

Generalized Lie derivations on the skew elements of a prime algebra with involution are characterised, provided the algebra does not satisfy a stated polynomial identity~\cite{WlGeneralizedLieDerivationsOnSke2011}.
